% GENERATED FILE. Edit canonical lessons, then run npm run generate:books.
% canonical-source-hash: fnv1a64:e1e1f2d674e91e01
% canonical-lessons: AR-C455-muntaj, AR-C455-marka, AR-C455-wajiha, AR-C455-arada-display, AR-C455-basta

\chapter{Shop windows}
\label{ch:ar-shop-windows}
\hlchaptermodality{455}

\begin{chapteropening}
\textbf{By the end of this chapter:} \emph{I can say a product, a brand, a shop window, to display and a market stall.}

Complete the last lesson of chapter 455: I can say a product, a brand, a shop window, to display and a market stall.
\end{chapteropening}

\section[\texorpdfstring{\ar{منتج} (muntaj)}{منتج (muntaj)}]{\ar{منتج} (muntaj) — a product}

\label{lesson:AR-C455-muntaj}



\begin{quote}
Before the new one: say the Arabic for a swimming pool, then the Arabic for a shopping centre.
\end{quote}

\subsection*{You'll want to know: \ar{منتج}}
\textbf{\ar{منتج}} — \emph{muntaj} — \textquotedblleft{}a product\textquotedblright{}.

\begin{quote}
\ar{هذا} \ar{منتج} \ar{محلي} \ar{جيد}.
\end{quote}

\begin{quote}
\emph{hādhā muntaj maḥallī jayyid.}
\end{quote}

\begin{quote}
This is a good local product.
\end{quote}

\begin{cousinweb}
\textbf{\ar{منتج}} (\emph{muntaj}) is built on the native Arabic root \emph{n-t-j}, \textquotedblleft{}to bring forth young\textquotedblright{} (Lane). It is the Form IV passive participle, \textquotedblleft{}what is brought forth\textquotedblright{}; \emph{natīja}, \textquotedblleft{}a result\textquotedblright{}, shares the root.
\end{cousinweb}

\begin{grammarlens}[title={What you've built}]
One more word for getting around.
\end{grammarlens}

\subsection*{Your turn}
\begin{itemize}
\raggedright
  \item \emph{Say it:} \emph{muntaj}
  \item \emph{Say it:} \emph{muntaj}, once more
  \item \emph{Say it:} \emph{hādhā muntaj maḥallī jayyid.} — This is a good local product.
  \item \emph{Recall:} say the Arabic for a swimming pool, then the Arabic for a shopping centre, then say \emph{muntaj} again
\end{itemize}

\begin{tcolorbox}[breakable,colback=teal!4,colframe=teal!35!black,title={Before you move on}]
What does \ar{منتج} mean? (\textquotedblleft{}A product\textquotedblright{}.) Say it once more.
\end{tcolorbox}
\section[\texorpdfstring{\ar{ماركة} (mārka)}{ماركة (mārka)}]{\ar{ماركة} (mārka) — a brand}

\label{lesson:AR-C455-marka}



\begin{quote}
Before the new one: say the Arabic for a shopping centre, then the Arabic for a product.
\end{quote}

\subsection*{You'll want to know: \ar{ماركة}}
\textbf{\ar{ماركة}} — \emph{mārka} — \textquotedblleft{}a brand\textquotedblright{}.

\begin{quote}
\ar{اشترى} \ar{أخي} \ar{حذاء} \ar{من} \ar{ماركة} \ar{مشهورة}.
\end{quote}

\begin{quote}
\emph{ishtarā akhī ḥidhāʾ min mārka mashhūra.}
\end{quote}

\begin{quote}
My brother bought shoes from a famous brand.
\end{quote}

Formal writing also says \emph{ʿalāma tijāriyya}.

\begin{cousinweb}
\textbf{\ar{ماركة}} (\emph{mārka}) is usually traced to Italian \emph{marca}, \textquotedblleft{}a mark\textquotedblright{}.
\end{cousinweb}

\begin{grammarlens}[title={What you've built}]
One more word for getting around.
\end{grammarlens}

\subsection*{Your turn}
\begin{itemize}
\raggedright
  \item \emph{Say it:} \emph{mārka}
  \item \emph{Say it:} \emph{mārka}, once more
  \item \emph{Say it:} \emph{ishtarā akhī ḥidhāʾ min mārka mashhūra.} — My brother bought shoes from a famous brand.
  \item \emph{Recall:} say the Arabic for a shopping centre, then the Arabic for a product, then say \emph{mārka} again
\end{itemize}

\begin{tcolorbox}[breakable,colback=teal!4,colframe=teal!35!black,title={Before you move on}]
What does \ar{ماركة} mean? (\textquotedblleft{}A brand\textquotedblright{}.) Say it once more.
\end{tcolorbox}
\section[\texorpdfstring{\ar{واجهة} (wājiha)}{واجهة (wājiha)}]{\ar{واجهة} (wājiha) — a shop window; a front}

\label{lesson:AR-C455-wajiha}



\begin{quote}
Before the new one: say the Arabic for a product, then the Arabic for a brand.
\end{quote}

\subsection*{You'll want to know: \ar{واجهة}}
\textbf{\ar{واجهة}} — \emph{wājiha} — \textquotedblleft{}a shop window; a front\textquotedblright{}.

\begin{quote}
\ar{رأيت} \ar{فستانا} \ar{جميلا} \ar{في} \ar{واجهة} \ar{المتجر}.
\end{quote}

\begin{quote}
\emph{raʾaytu fustānan jamīlan fī wājihat al-matjar.}
\end{quote}

\begin{quote}
I saw a beautiful dress in the shop window.
\end{quote}

\begin{cousinweb}
\textbf{\ar{واجهة}} (\emph{wājiha}) is built on the native Arabic root \emph{w-j-h}, \textquotedblleft{}the face; to face\textquotedblright{} (Lane). It is the part of a building that faces the street.
\end{cousinweb}

\begin{grammarlens}[title={What you've built}]
One more word for getting around.
\end{grammarlens}

\subsection*{Your turn}
\begin{itemize}
\raggedright
  \item \emph{Say it:} \emph{wājiha}
  \item \emph{Say it:} \emph{wājiha}, once more
  \item \emph{Say it:} \emph{raʾaytu fustānan jamīlan fī wājihat al-matjar.} — I saw a beautiful dress in the shop window.
  \item \emph{Recall:} say the Arabic for a product, then the Arabic for a brand, then say \emph{wājiha} again
\end{itemize}

\begin{tcolorbox}[breakable,colback=teal!4,colframe=teal!35!black,title={Before you move on}]
What does \ar{واجهة} mean? (\textquotedblleft{}A shop window; a front\textquotedblright{}.) Say it once more.
\end{tcolorbox}
\section[\texorpdfstring{\ar{عرض} (ʿaraḍa)}{عرض (ʿaraḍa)}]{\ar{عرض} (ʿaraḍa) — to display; to offer}

\label{lesson:AR-C455-arada-display}



\begin{quote}
Before the new one: say the Arabic for a brand, then the Arabic for a shop window.
\end{quote}

\subsection*{You'll want to know: \ar{عرض}}
\textbf{\ar{عرض}} — \emph{ʿaraḍa} — \textquotedblleft{}to display; to offer\textquotedblright{}.

\begin{quote}
\ar{عرض} \ar{المتجر} \ar{ملابس} \ar{الشتاء} \ar{في} \ar{الواجهة}.
\end{quote}

\begin{quote}
\emph{ʿaraḍa al-matjar malābis ash-shitāʾ fī al-wājiha.}
\end{quote}

\begin{quote}
The shop displayed winter clothes in the window.
\end{quote}

Present: \ar{يعرض}, \emph{yaʿriḍ}.

The noun \emph{ʿarḍ} is an offer or a show.

\begin{cousinweb}
\textbf{\ar{عرض}} (\emph{ʿaraḍa}) is built on the native Arabic root \emph{ʿ-r-ḍ}, \textquotedblleft{}to show, to present\textquotedblright{} (Lane). It is the plain verb; \emph{ʿarīḍ}, \textquotedblleft{}wide\textquotedblright{}, shares the root.
\end{cousinweb}

\begin{grammarlens}[title={What you've built}]
One more word for getting around.
\end{grammarlens}

\subsection*{Your turn}
\begin{itemize}
\raggedright
  \item \emph{Say it:} \emph{ʿaraḍa}
  \item \emph{Say it:} \emph{ʿaraḍa}, once more
  \item \emph{Say it:} \emph{ʿaraḍa al-matjar malābis ash-shitāʾ fī al-wājiha.} — The shop displayed winter clothes in the window.
  \item \emph{Recall:} say the Arabic for a brand, then the Arabic for a shop window, then say \emph{ʿaraḍa} again
\end{itemize}

\begin{tcolorbox}[breakable,colback=teal!4,colframe=teal!35!black,title={Before you move on}]
What does \ar{عرض} mean? (\textquotedblleft{}To display; to offer\textquotedblright{}.) Say it once more.
\end{tcolorbox}
\section[\texorpdfstring{\ar{بسطة} (basṭa)}{بسطة (basṭa)}]{\ar{بسطة} (basṭa) — a market stall}

\label{lesson:AR-C455-basta}



\begin{quote}
Before the new one: say the Arabic for a shop window, then the Arabic for to display.
\end{quote}

\subsection*{You'll want to know: \ar{بسطة}}
\textbf{\ar{بسطة}} — \emph{basṭa} — \textquotedblleft{}a market stall\textquotedblright{}.

\begin{quote}
\ar{يبيع} \ar{جاري} \ar{الفاكهة} \ar{على} \ar{بسطة} \ar{في} \ar{السوق}.
\end{quote}

\begin{quote}
\emph{yabīʿu jārī al-fākiha ʿalā basṭa fī as-sūq.}
\end{quote}

\begin{quote}
My neighbour sells fruit on a stall in the market.
\end{quote}

The word is common in the Levant.

\begin{cousinweb}
\textbf{\ar{بسطة}} (\emph{basṭa}) is built on the native Arabic root \emph{b-s-ṭ}, \textquotedblleft{}to spread out\textquotedblright{} (Lane). A \emph{basṭa} is goods spread out for sale.
\end{cousinweb}

\begin{grammarlens}[title={What you've built}]
One more word for getting around.
\end{grammarlens}

\subsection*{Your turn}
\begin{itemize}
\raggedright
  \item \emph{Say it:} \emph{basṭa}
  \item \emph{Say it:} \emph{basṭa}, once more
  \item \emph{Say it:} \emph{yabīʿu jārī al-fākiha ʿalā basṭa fī as-sūq.} — My neighbour sells fruit on a stall in the market.
  \item \emph{Recall:} say the Arabic for a shop window, then the Arabic for to display, then say \emph{basṭa} again
\end{itemize}

\begin{tcolorbox}[breakable,colback=teal!4,colframe=teal!35!black,title={Before you move on}]
What does \ar{بسطة} mean? (\textquotedblleft{}A market stall\textquotedblright{}.) Say it once more.
\end{tcolorbox}
